\documentclass{situationist-preprint}

\title{Kierkegaard's Leap and Evidential Norms}
\msid{REL-17-b}
\msfamily{REL-17}
\msdiscipline{Philosophy of Religion}
\mstype{Consequence}
\mskeywords{Kierkegaard; leap of faith; objective uncertainty; evidentialism; Clifford; James}

\begin{document}
\maketitle

\begin{abstract}
Kierkegaard's pseudonym Johannes Climacus holds that faith is a leap, made in objective uncertainty and with infinite passion, and that it holds fast to what the understanding finds absurd. He is often read as defying evidential norms. We derive from the conclusion of REL-17-a a more exact verdict. Where the absurd is read as what the understanding cannot settle, the leap satisfies the conditions under which faith is rational, and evidentialism neither requires nor forbids it. Where the absurd is read as what the evidence tells against, evidentialism excludes it. Clifford's principle, on the view of EPI-23, does not bind the leap in either case.
\end{abstract}

\section{Three claims}

In the \emph{Concluding Unscientific Postscript}, Climacus makes three claims about faith \citep{Kierkegaard1846}. (K1) Objective inquiry, historical or speculative, is an approximation that never yields a decision about the matters on which faith turns. (K2) The believer's interest in those matters is infinite: they concern her eternal happiness, and so how she lives. (K3) Faith holds fast to the absolute paradox, which is absurd to the understanding.

REL-17-a defines faith that $p$ as a practical commitment to $p$ together with the absence of disbelief (Definition 2.3), and concludes that faith is rational where the evidence underdetermines $p$ and the stakes are practical (Theorem 3.4). Evidentialism assesses doxastic attitudes only (Principle P1); it excludes faith only where the evidence opposes $p$ (Corollary 3.5).

\begin{proposition}\label{p:leap}
On the reading of (K3) on which the absurd is what the understanding cannot settle, the leap is rational, and evidentialism neither requires nor forbids it.
\end{proposition}

\begin{proof}
(K1) says that the evidence leaves the question underdetermined in the sense of Definition 2.1 of REL-17-a, and that further inquiry will not change this. (K2) says that the agent faces a choice between ways of life whose value depends on the answer, and, for one who makes the leap, that commitment has by her lights at least the value of abstention; the stakes are practical in the sense of Definition 2.4. The leap is the adoption of a commitment, not of a belief fitted to evidence: Climacus insists that faith is not the result of an approximation. On this reading (K3) adds only that the understanding does not settle the matter. By Theorem 3.4 and Propositions 3.2 and 3.3 of REL-17-a, the leap is rational and evidentialism neither forbids nor requires it.
\end{proof}

\begin{proposition}\label{p:absurd}
On the reading of (K3) on which the absurd is what the evidence tells against, the leap is excluded by evidentialism.
\end{proposition}

\begin{proof}
If the evidence opposes $p$, disbelief is the only fitting attitude (Definition 2.2 of REL-17-a), and faith requires its absence. By Corollary 3.5 of REL-17-a, faith that $p$ is then irrational.
\end{proof}

The textual case favors the first reading: Climacus describes the paradox as what the understanding cannot think through, and he is as hostile to the claim that the evidence refutes Christianity as to the claim that it proves it. \citet{Evans1998} defends a reading of this kind, on which faith is above reason rather than against it. Proposition~\ref{p:absurd} shows that the second reading cannot be made rational by appeal to passion alone.

\section{Clifford and James}

\citet{Clifford1877} held that it is wrong always and for anyone to believe on insufficient evidence; \citet{James1896} defended the right to believe on a genuine option that the intellect cannot decide. EPI-23 concludes that Clifford's principle holds where the stakes are high and evidence is available, and that James's right to believe applies where neither condition holds.

\begin{objection}
The stakes of the leap are infinitely high. If any belief falls under Clifford's principle, this one does.
\end{objection}

\begin{reply}
Clifford's principle, on the view of EPI-23, binds where the stakes are high \emph{and} evidence is available. (K1) says precisely that evidence able to decide the matter is not available, and that inquiry does not make it so. The second condition fails, and the principle does not apply. More fundamentally, Clifford's principle is a norm on belief, and the leap, on Proposition~\ref{p:leap}, is a commitment. Climacus's own contrast between faith and the results of approximation supports this: what he requires of the believer is a resolution in how to live, not an assent proportioned to grounds.
\end{reply}

\section{Conclusion}

Kierkegaard's leap is not a defiance of evidential norms. Read as a commitment made where inquiry cannot decide and life cannot wait, it is exactly the kind of faith that REL-17-a shows to be rational, and evidentialism neither requires nor forbids it. Read as an embrace of what the evidence refutes, it is irrational, and no intensity of passion alters that.

\begin{thebibliography}{99}\small
\bibitem[Clifford(1877)]{Clifford1877} Clifford, W.~K. (1877). The ethics of belief. \emph{Contemporary Review} 29.
\bibitem[Evans(1998)]{Evans1998} Evans, C.~S. (1998). \emph{Faith Beyond Reason: A Kierkegaardian Account}. Eerdmans.
\bibitem[James(1896)]{James1896} James, W. (1896). The will to believe. \emph{The New World} 5.
\bibitem[Kierkegaard(1846)]{Kierkegaard1846} Kierkegaard, S. (1846). \emph{Concluding Unscientific Postscript to Philosophical Fragments}. Trans. H.~V. Hong and E.~H. Hong, Princeton University Press, 1992.
\end{thebibliography}

\end{document}
